
%%Mẫu ĐA/KLTN TRƯỜNG ĐẠI HỌC CÔNG NGHỆ THÔNG TIN VÀ TRUYỀN THÔNG (ICTU)

%%%%%%%%%%%%%%%%%%%%%%%%%%
\documentclass[12pt,a4paper, reqno, oneside]{reportICTU}
%\usepackage[13pt]{vnextsizes}
\usepackage{amsmath,amsxtra,amssymb,latexsym, amscd,amsthm,multirow} %gói công thức toán
\usepackage{enumerate} 
\usepackage{graphicx,subfigure,psfrag} 
\usepackage[top=2.0cm, bottom=2.0cm, left=3.0cm, right=2.0cm, headheight=15pt]{geometry} % căn lề
\usepackage{indentfirst}
\usepackage{color}
\usepackage[utf8]{vietnam} % Tiếng Việt 
 
\usepackage[mathscr]{eucal}
\usepackage{amsfonts}
\usepackage{tikz}
\usetikzlibrary{calc}
\usepackage{setspace} % Bắt buộc phải có để sử dụng môi trường spacing
%\usepackage{xhfill}
%{\setstretch{1.2}

\usepackage{fancybox} 
\usepackage{boxedminipage}
\usepackage{makecell}
\newlength{\dll}
\newenvironment{mdbox}
{\setlength{\fboxsep}{10pt}\setlength{\shadowsize}{0pt}\setlength{\dll}{\linewidth}
\addtolength{\dll}{-2\fboxsep}
\addtolength{\dll}{-2\fboxrule}
\begin{Sbox}\begin{minipage}{\dll}}
{\end{minipage}\end{Sbox}
\noindent\shadowbox{\TheSbox}}

\usepackage[displaymath, mathlines]{lineno}%đánh số dòng


\usepackage{algorithm}
\usepackage{algpseudocode}

\floatname{algorithm}{Thuật toán}
\renewcommand{\algorithmicrequire}{\textbf{Đầu vào:}}
\renewcommand{\algorithmicensure}{\textbf{Đầu ra:}}


\DeclareMathOperator*{\argmax}{arg\,max}


% \usepackage[subfigure]{tocloft}
\usepackage{setspace}

\usepackage{scrextend}
\changefontsizes[16pt]{13pt} %font 13

\usepackage{caption} % caption table and figure
\captionsetup{margin=10pt,font=small}%,labelfont=bf}


\usepackage{sansmathaccent}
\pdfmapfile{+sansmathaccent.map}
%\input setbmp
\usepackage[unicode]{hyperref}
\usepackage{titledot}
\titlespacing{\chapter}{0pt}{-10pt}{0pt}[0pt]
\titlename{Chương}
%+ Can le van ban

%%% Định nghĩa mỗi trường 
%\newtheorem*{remarks}{Remarks}
\newtheorem{lemma}{\bf Bổ đề}[section]
\newtheorem{dl}{\bf Định lý}[section]
\newtheorem{theorem}{\bf Định lý}[section]
\newtheorem{proposition}{\bf Mệnh đề}[section]
\newtheorem{corollary}{\bf Hệ quả}[section]
\newtheorem{definition}{\bf Định nghĩa}[section]
\newtheorem{remark}{\bf Chú ý}[section]
\newtheorem{NX}{\bf Nhận xét}[section]
\newtheorem{Bang}{\bf Bảng}[section]
\newtheorem{vd}{\bf Ví dụ}[section]
\newtheorem{example}{Ví dụ}[section]
\newtheorem{Hi}[theorem]{\bf Hình}

\newcommand{\pf}{\noindent{\bf Proof.} }
\newcommand{\eop}{{ \vrule height7pt width7pt depth0pt}\par\bigskip}
\newcommand{\meop}{\quad {\vrule height7pt width7pt depth0pt}}



\def\figpath{Figures/} %thư mục chứa hình vẽ

 
%% Định nghĩa một số ký hiệu toán học
\def\RR{\mathbb{R}}
\DeclareMathOperator{\sep}{sep}
\DeclareMathOperator{\dist}{dist}
\def\eps{\varepsilon}

% \setlength{\oddsidemargin}{0in}
% \setlength{\textwidth}{6.truein}
% \setlength{\topmargin}{-0.5truein}
% \setlength{\textheight}{9.3truein}
%\renewcommand{\baselinestretch}{1.5}
\usepackage{mathptmx}
%+ Dinh nghia mot so ky hieu toan hoc.


\newtheorem{prof}{}
%+ Bat dau tai lieu tai day.
 %\usepackage[usehighlevels]{alnumsec}
\usepackage{fancyhdr}
\pagestyle{fancy}
\rhead{}
\chead{\thepage}
\lhead{}
\lfoot{}
\cfoot{}
\rfoot{}
\renewcommand{\headrulewidth}{0pt}
\renewcommand{\footrulewidth}{0pt}
%\newtheorem{algorithm}{Thuật toán}
%\newtheorem{testproblem}{Test Problem}
\newtheorem{testproblem}{Bài toán}



\begin{document}


% \alnumsecstyle{Raaa}
% \surroundRoman{}{}
% \surroundarabic[(][)]{.}{}

%+ ==========================
\newpage
\begin{titlepage}

\begin{tikzpicture}[remember picture,overlay]
    \draw [line width=1pt]
    ($ (current page.north west) + (3cm,-2.5cm) $)
    rectangle
    ($ (current page.south east) + (-1.5cm,2.5cm) $);
\draw [line width=3pt]
    ($ (current page.north west) + (3.15cm,-2.65cm) $)
    rectangle
    ($ (current page.south east) + (-1.65cm,2.65cm) $); 
\end{tikzpicture}

\begin{center}

\begin{spacing}{1}
\textbf{TRƯỜNG ĐẠI HỌC CÔNG NGHỆ THÔNG TIN VÀ TRUYỀN THÔNG} \\
\textbf{KHOA \ldots\\[1cm]}
\end{spacing}

\includegraphics[scale=0.04]{\figpath logo_ictu.jpg}\\[1cm]

\begin{large}
\begin{spacing}{1}
\textbf{ĐỒ ÁN/KHÓA LUẬN TỐT NGHIỆP \\[2.5cm]}
\end{spacing}
\end{large}

\textbf{ĐỀ TÀI:} \\[0.3cm]
\textbf{{\Large PHƯƠNG PHÁP RBF-FD ...}} \\[2.5cm]

\begin{tabular}{ r l }
\textbf{Giảng viên hướng dẫn:} & \textbf{TS. Nguyễn Văn ....} \\
\textbf{Sinh viên thực hiện:} & \textbf{Trần Thị Thu ....} \\
\textbf{Lớp:} & \textbf{Công nghệ thông tin K22A} \\
\textbf{Mã sinh viên:} & \textbf{DTC235123456} \\
\end{tabular}
\\[6cm]

\textbf{Thái Nguyên, năm 2026}

\end{center}

\end{titlepage}





\begin{titlepage}

\begin{tikzpicture}[remember picture,overlay]
    \draw [line width=1pt]
    ($ (current page.north west) + (3cm,-2.5cm) $)
    rectangle
    ($ (current page.south east) + (-1.5cm,2.5cm) $);
\draw [line width=3pt]
    ($ (current page.north west) + (3.15cm,-2.65cm) $)
    rectangle
    ($ (current page.south east) + (-1.65cm,2.65cm) $); 
\end{tikzpicture}

\begin{center}

\begin{spacing}{1}
\textbf{TRƯỜNG ĐẠI HỌC CÔNG NGHỆ THÔNG TIN VÀ TRUYỀN THÔNG} \\
\textbf{KHOA .....\\[1.5cm]}
\end{spacing}

\includegraphics[scale=0.04]{\figpath logo_ictu.jpg}\\[1cm]

\begin{large}
\begin{spacing}{1}
\textbf{ĐỒ ÁN/KHÓA LUẬN TỐT NGHIỆP} \\[2cm]
\end{spacing}
\end{large}

\textbf{ĐỀ TÀI:} \\[0.3cm]
\textbf{{\Large PHƯƠNG PHÁP RBF-FD ....}} \\[2cm]

\begin{tabular}{ r l }
\textbf{Giảng viên hướng dẫn:} & \textbf{TS. Nguyễn Văn A} \\
\textbf{Sinh viên thực hiện:} & \textbf{Trần Thị B} \\
\textbf{Lớp:} & \textbf{Công nghệ thông tin K2sA} \\
\textbf{Mã sinh viên:} & \textbf{DTC23123456} \\
\end{tabular}
\\[6cm]

\textbf{Thái Nguyên, năm 2026}

\end{center}

\end{titlepage}





%+ Page style
\setlength{\baselineskip}{22truept}
\pagenumbering{roman}
{
%\newgeometry{top=3.0cm, bottom=2.0cm, left=3cm, right=2.0cm}
\begin{spacing}{1.3}
\addcontentsline{toc}{chapter}{LỜI CAM ĐOAN}
\vspace*{5pt}


\noindent
\begin{minipage}[t]{0.45\textwidth}
    \centering
            ĐẠI HỌC THÁI NGUYÊN \\
    \textbf{TRƯỜNG ĐẠI HỌC CÔNG NGHỆ} \\
    \textbf{THÔNG TIN VÀ TRUYỀN THÔNG}
\end{minipage}%
\begin{minipage}[t]{0.59\textwidth}
    \centering
    \textbf{CỘNG HÒA XÃ HỘI CHỦ NGHĨA VIỆT NAM} \\
    \textbf{Độc lập - Tự do - Hạnh phúc} \\
    \setlength{\baselineskip}{10truept}
    \rule{0.4\linewidth}{0.4pt}
\end{minipage}

\vspace{1cm}

\begin{center}
    \textbf{\large LỜI CAM ĐOAN}\\[0.3cm]
    \textit{(kèm theo Thông báo số: \ldots{} /TB-ĐH CNTT\&TT ngày \ldots{} tháng \ldots\ldots{} năm 2026}\\
    \textit{của Hiệu trưởng Trường ĐH CNTT\&TT)}\\[0.1cm]
    \setlength{\baselineskip}{10truept}
    \rule{0.4\textwidth}{0.4pt}
\end{center}

\vspace{0.5cm}

{
\setlength{\baselineskip}{28truept}
\hspace{1cm} Em xin cam đoan đã thực hiện việc kiểm tra mức độ tương đồng nội dung (ĐA/KLTN) qua phần mềm Turnitin một cách trung thực và đạt kết quả mức độ tương đồng \ldots\ldots\ldots\ldots{} \%. Bản đề tài (ĐA/KLTN) kiểm tra qua phần mềm là bản cứng đã nộp để bảo vệ trước hội đồng.

\hspace{1cm} Nếu sai em xin hoàn toàn chịu trách nhiệm.

\hspace{1cm} Xin chân thành cảm ơn./.
}

\vspace{1cm}

\begin{tabular}{p{0.35\textwidth} >{\centering\arraybackslash}p{0.55\textwidth}}
 & \textit{Thái Nguyên, ngày \ldots{} tháng \ldots{} năm 202\ldots} \\[0.5cm]
 & \textbf{TÁC GIẢ CỦA SẢN PHẨM} \\
 & \textit{(Ký và ghi rõ họ tên)} \\[2.5cm]
\end{tabular}
 % LỜI CAM ĐOAN
\newpage
\addcontentsline{toc}{chapter}{LỜI CẢM ƠN}
\vspace*{15pt}
\begin{center}
\large{\textbf{LỜI CẢM ƠN}}
\end{center}



Sinh viên trình bày lời cảm ơn tới nhà trường, khoa, giảng viên hướng dẫn, gia đình, bạn bè và các cá nhân, tổ chức đã hỗ trợ trong quá trình thực hiện khóa luận.

\vspace{1cm}

\begin{flushright}
Thái Nguyên, năm \ldots

Sinh viên

\vspace{2cm}

Trần Thị B
\end{flushright}


  %{LỜI CẢM ƠN
\newpage
\addcontentsline{toc}{chapter}{DANH MỤC TỪ VIẾT TẮT}
\vspace*{15pt}
\begin{center}
\large{\textbf{DANH MỤC TỪ VIẾT TẮT}}
\end{center}


\begin{table}[H]
\centering
\renewcommand{\arraystretch}{1.3}
\begin{tabular}{|p{4cm}|p{9cm}|}
\hline
\textbf{Từ viết tắt} & \textbf{Diễn giải} \\
\hline
CNTT & Công nghệ thông tin \\
\hline
API & Application Programming Interface \\
\hline
UI & User Interface \\
\hline
\end{tabular}
\end{table}


 %DANH MỤC TỪ VIẾT TẮT

\newpage
\tableofcontents
\newpage
\addcontentsline{toc}{chapter}{Danh sách bảng}%section
\listoftables
\newpage
\addcontentsline{toc}{chapter}{Danh sách hình vẽ}%section
\listoffigures
\end{spacing}
%\newpage

}

%%% ==== NỘI DUNG CHÍNH
\newpage
\pagenumbering{arabic}
{

\addcontentsline{toc}{chapter}{PHẦN MỞ ĐẦU}
\vspace*{15pt}
\begin{center}
\large{\textbf{PHẦN MỞ ĐẦU}}
\end{center}

\section*{Lý do chọn đề tài}
Sinh viên trình bày bối cảnh, vấn đề thực tế và lý do cần thực hiện đề tài.

\section*{Mục tiêu nghiên cứu}
Sinh viên nêu mục tiêu tổng quát và các mục tiêu cụ thể của đề tài.

\section*{Đối tượng và phạm vi nghiên cứu}
Sinh viên mô tả đối tượng nghiên cứu, phạm vi thực hiện, dữ liệu sử dụng và các giới hạn của đề tài.

\section*{Phương pháp nghiên cứu}
Sinh viên trình bày phương pháp khảo sát, phân tích, thiết kế, triển khai, kiểm thử và đánh giá.

\section*{Cấu trúc báo cáo}
Báo cáo gồm phần mở đầu, các chương nội dung chính, kết luận và tài liệu tham khảo.


 % PHẦN MỞ ĐẦU
 \newpage
 \setcounter{chapter}{0}
\setcounter{section}{0}
\setcounter{table}{0}
\setcounter{figure}{0}
%\addcontentsline{toc}{chapter}{CƠ SỞ LÝ THUYẾT}
\chapter{CƠ SỞ LÝ THUYẾT}

\textit{Chương này trình bày những công nghệ chính sử dụng trong đồ án.}

\section{Tổng quan đề tài}
Sinh viên trình bày tổng quan về đề tài, bối cảnh ứng dụng và vấn đề cần giải quyết.

\section{Cơ sở lý thuyết}
Sinh viên trình bày các khái niệm, mô hình, thuật toán, phương pháp hoặc công nghệ nền tảng liên quan đến đề tài.
Ví dụ, khi trình bày tài liệu nền tảng, sinh viên có thể trích dẫn như sau: quy trình phát triển phần mềm được trình bày trong tài liệu của Sommerville~\cite{sommerville2016software}.

Ví dụ danh sách gạch đầu dòng:
\begin{itemize}
    \item Công nghệ hoặc nền tảng thứ nhất sử dụng trong đồ án.
    \item Công nghệ hoặc nền tảng thứ hai sử dụng trong đồ án.
    \item Công nghệ hoặc nền tảng thứ ba sử dụng trong đồ án.
\end{itemize}

\section{Các công trình hoặc hệ thống liên quan}
Sinh viên tóm tắt các nghiên cứu, sản phẩm hoặc hệ thống tương tự. Nên nêu ưu điểm, hạn chế và bài học kế thừa.

\section{Định hướng giải quyết}
Sinh viên nêu hướng tiếp cận dự kiến, các giả định ban đầu và lý do lựa chọn phương án thực hiện.


  % Nội dung chương 1
 \setcounter{chapter}{1}
\setcounter{chapter}{1} 
\setcounter{table}{0}
\setcounter{figure}{0}
%\addcontentsline{toc}{chapter}{KHẢO SÁT VÀ PHÂN TÍCH HỆ THỐNG}
\chapter{KHẢO SÁT VÀ PHÂN TÍCH HỆ THỐNG}


\section{Khảo sát hiện trạng}
Sinh viên mô tả hiện trạng, quy trình nghiệp vụ hoặc bài toán thực tế liên quan đến đề tài.

\section{Phân tích yêu cầu}
Sinh viên trình bày yêu cầu chức năng, yêu cầu phi chức năng và các ràng buộc của hệ thống hoặc giải pháp.

\section{Phương pháp thực hiện}
Sinh viên mô tả phương pháp, quy trình, công cụ và dữ liệu sử dụng trong quá trình thực hiện đề tài.

\section{Kế hoạch thực hiện}
Sinh viên trình bày các giai đoạn công việc, mốc thời gian và sản phẩm dự kiến của từng giai đoạn.
Ví dụ, kế hoạch thực hiện có thể được trình bày như Bảng~\ref{tab:ke-hoach}.

\begin{table}[H]
\centering
\caption{Ví dụ bảng kế hoạch thực hiện}
\label{tab:ke-hoach}
\begin{tabular}{|c|p{5cm}|p{5cm}|}
\hline
\textbf{STT} & \textbf{Công việc} & \textbf{Kết quả dự kiến} \\
\hline
1 & Khảo sát và thu thập yêu cầu & Danh sách yêu cầu \\
\hline
2 & Phân tích và thiết kế & Tài liệu thiết kế \\
\hline
3 & Triển khai và kiểm thử & Sản phẩm và kết quả kiểm thử \\
\hline
\end{tabular}
\end{table}



\begin{figure}[htbp!]
 \begin{center}
\psfrag{z}{$\zeta$}
\psfrag{x}{$\xi$}
\psfrag{m}{$\xi_{\rm mid}$}
\psfrag{p}{$\xi_{\rm mid}^+$}
\psfrag{mm}{$\xi_{\rm mid}^-$}
\psfrag{xp}{$\xi_+$}
\psfrag{xm}{$\xi_-$}
\psfrag{xps}{$\xi'_+$}
\psfrag{xms}{$\xi'_-$}
\includegraphics[width=5.5cm]{\figpath refine.eps} 
\end{center}
\caption{Các tâm được thêm vào tại lân cận cạnh được đánh dấu. 
}
\label{refine}
\end{figure}
 Thuật toán trong \cite{DavyOanh11} làm mịn cạnh $\zeta \xi$ bằng cách chèn thêm tâm mới $\xi _ { \rm {mid } }$ là trung điểm của cạnh này. Còn trong \cite{ODP2017}, ngoài trung điểm còn thêm các điểm $ \xi _ { \rm{ mid } } ^ { + } , \xi _ { \rm{ mid } } ^ { - }$ nằm trên đường trung của cạnh $\zeta \xi$ và đối xứng nhau qua $\xi _ { \rm{mid }}$. Tuy nhiên các điểm này có thể nằm ở vị trí rất gần điểm đã tồn tại hoặc một điểm mới được tạo ra khi làm mịn một cạnh khác. Để tránh điều này các thuật toán trong  \cite{DavyOanh11, ODP2017} chỉ làm mịn nếu 
\begin{equation}\label{mu}
\dist(\xi',\Xi')\ge\mu\sep_{\xi'}(\Xi'),\quad\text{với $\mu\in(0,1)$,}
\end{equation}
tức là, việc chèn một tâm mới $\xi'$ trong tập hiện tại $\Xi'$ không làm giảm  
\emph{khoảng cách tách biệt địa phương} được xác định công thức
\begin{equation}\label{sep}
\sep_{\xi'}(\Xi'):=\frac{1}{4}\sum_{i=1}^4\dist(\xi_i,\Xi'\setminus\{\xi_i\}),
\end{equation}
trong đó $\xi_1,\ldots,\xi_4$ là bốn điểm gần nhất trong $\Xi'$  đến $\xi'$, và
$\dist(x,Y):=\inf\{\|x-y\|:\,y\in Y\}$
là khoảng cách từ điểm $x$ đến tập $Y$. 
% Tuy nhiên, chúng tôi đã kiểm tra 
% $\xi'\notin\partial\Omega$ không ở vị trí quá gần biên, 
% xem các Bước II.2 và II.3.i của Thuật toán~\ref{alg2Refi}.  

Nếu $\xi \in \partial \Xi$, thì tìm 2 điểm lân cận $\xi _ { - } , \xi _ { + } \in \partial \Xi$ và chèn thêm 2 trung điểm $\xi _ { - } ^ { \prime } , \xi _ { + } ^ { \prime }$ của các cạnh $\xi \xi _ { - }$ và $\xi \xi _ { + }$.



% 

\begin{algorithm}\label{alg1Refi}\rm {\bf Thuật toán làm mịn thích nghi  [Thuật toán 2, \cite{DavyOanh11}]} \\
\emph{Input:} Bộ tâm $\Xi$ và các khuôn điểm $\{\Xi_\zeta:\,\zeta\in\Xi_{int}\}$. \\
\emph{Output:} Bộ tâm được làm mịn $\Xi'$ và các giá stencil $\{ \Xi' _ { \zeta }: \zeta \in \Xi' \setminus \partial \Xi' \}$ . \\
\emph{Các tham số:}   $\gamma = 0.5$, $\mu = 0.7$. 
\begin{enumerate}[I.] 
\item\label{1step1} Tính độ lệch lớn nhất  $\bar\eps=\bar\eps(\Xi)$ và đánh dấu tất cả các cạnh $\zeta\xi, \xi \in \Xi_{\zeta}\setminus \{\zeta\}, \zeta \in \Xi_{\rm int}$ sao cho $\eps(\zeta, \xi) \ge \gamma\bar\eps$. Thiết lập  $\Xi':=\Xi$.
\item \label{1step2} \emph{Với mỗi} cạnh được đánh dấu $\zeta\xi$: 
\begin{enumerate}[1.]
\item Đặt $\xi':=(\zeta+\xi)/2$.
\item \emph{Nếu} $\dist\left( \xi' , \Xi' \right) \ge \mu \sep _ { \xi' } \left( \Xi' \right)$:
\begin{enumerate}[i.] 
\item Thiết lập  $\Xi': = \Xi' \cup \{ \xi'\}$.
\item \emph{Nếu} $\xi \in \partial \Xi$ thì thiết lập $\Xi': = \Xi' \cup \{ \xi' _ { - }, \xi' _ { + }\}$.
\end{enumerate}
\end{enumerate}
\item 
\begin{enumerate}[1.]
\item \emph{Với mỗi} $\zeta \in \Xi' \setminus \partial \Xi'$:
\begin{enumerate}[i.] 
\item Áp dụng Thuật toán \ref{alg1} để tìm $\Xi' _ { \zeta } \subset \Xi'$. 
\item Tính khoảng cách $d _ { \xi } = \| \zeta - \xi \| _ { 2 } , \xi \in \Xi' _ { \zeta } \setminus \{ \zeta \}$ và trung bình của chúng 
$d _ { a v } : = \frac { 1 } { \ell } \sum _ { \xi \in \Xi' _ { \zeta } \setminus \{ \zeta \} } d _ { \xi }$, trong đó $\ell = \# \Xi' _ { \zeta } - 1$ và 
đánh dấu tất cả các cạnh $\zeta \xi , \xi \in \Xi _ { \zeta } \setminus \{ \zeta \}$ sao cho $d _ { \xi } > 2 d _ { a v }$.
\end{enumerate}
\item Lặp lại Bước II đối với cạnh đã được đánh dấu trong Bước III.1 và tập tất cả các tâm mới đã được tạo trong tập $\Xi_{aux}$.
\item Áp dụng Thuật toán \ref{alg1} để tìm $\Xi' _ { \zeta }$ với mọi $\zeta \in \Xi_{aux}$ và tính lại $\Xi' _ { \zeta }$ với $\zeta \in \Xi' \setminus \Xi _ { a u x }$ sao cho $\dist  \left( \zeta , \Xi _ { a u x } \right) < d \left( \Xi' _ { \zeta }\right): = \max _ { \xi \in \Xi'_ { \zeta } } \| \zeta - \xi \| _ { 2 }$. 
\end{enumerate}
\end{enumerate}
\end{algorithm}
%   % Nội dung chương 2
 \setcounter{chapter}{2}
\setcounter{chapter}{2} 
\setcounter{table}{0}
\setcounter{figure}{0}
%\addcontentsline{toc}{chapter}{KHẢO SÁT VÀ PHÂN TÍCH HỆ THỐNG}
\chapter{THIẾT KẾ HỆ THỐNG}

\section{Thiết kế tổng thể}
Sinh viên trình bày kiến trúc tổng thể, các thành phần chính và mối quan hệ giữa các thành phần.

\section{Thiết kế dữ liệu}
Sinh viên mô tả dữ liệu đầu vào, dữ liệu đầu ra, mô hình cơ sở dữ liệu hoặc cấu trúc lưu trữ.

\section{Thiết kế chức năng}
Sinh viên trình bày các chức năng chính, luồng xử lý và các biểu đồ minh họa nếu có.

\section{Thiết kế giao diện hoặc mô hình thực nghiệm}
Sinh viên trình bày bản thiết kế giao diện, mô hình thí nghiệm, kịch bản sử dụng hoặc quy trình vận hành.
Ví dụ, vị trí chèn hình minh họa được thể hiện trong Hình~\ref{fig:vi-du-hinh}.

\begin{figure}[H]
\centering
\includegraphics[width=0.35\textwidth]{\figpath logo_ictu.jpg}
\caption{Ví dụ chèn ảnh trong báo cáo}
\label{fig:vi-du-hinh}
\end{figure}


 % Nội dung chương 3
 %\setcounter{chapter}{2}
\setcounter{chapter}{2} 
\setcounter{table}{0}
\setcounter{figure}{0}
%\addcontentsline{toc}{chapter}{KHẢO SÁT VÀ PHÂN TÍCH HỆ THỐNG}
\chapter{THIẾT KẾ HỆ THỐNG}

\section{Thiết kế tổng thể}
Sinh viên trình bày kiến trúc tổng thể, các thành phần chính và mối quan hệ giữa các thành phần.

\section{Thiết kế dữ liệu}
Sinh viên mô tả dữ liệu đầu vào, dữ liệu đầu ra, mô hình cơ sở dữ liệu hoặc cấu trúc lưu trữ.

\section{Thiết kế chức năng}
Sinh viên trình bày các chức năng chính, luồng xử lý và các biểu đồ minh họa nếu có.

\section{Thiết kế giao diện hoặc mô hình thực nghiệm}
Sinh viên trình bày bản thiết kế giao diện, mô hình thí nghiệm, kịch bản sử dụng hoặc quy trình vận hành.
Ví dụ, vị trí chèn hình minh họa được thể hiện trong Hình~\ref{fig:vi-du-hinh}.

\begin{figure}[H]
\centering
\includegraphics[width=0.35\textwidth]{\figpath logo_ictu.jpg}
\caption{Ví dụ chèn ảnh trong báo cáo}
\label{fig:vi-du-hinh}
\end{figure}


 % Nội dung chương 3
 \newpage
 \addcontentsline{toc}{chapter}{KẾT LUẬN VÀ KIẾN NGHỊ}
\begin{center}
\large {\textbf{KẾT LUẬN}}
\end{center}

\section*{Kết quả đạt được}
Sinh viên tóm tắt các kết quả chính đã thực hiện được so với mục tiêu ban đầu.

\section*{Hạn chế}
Sinh viên trình bày các hạn chế về phạm vi, dữ liệu, phương pháp, kỹ thuật hoặc kết quả thực nghiệm.

\section*{Hướng phát triển}
Sinh viên đề xuất các hướng cải tiến, mở rộng hoặc nghiên cứu tiếp theo.




}

%+ Tai lieu tham khao
 \newpage
 \setlength{\baselineskip}{19truept}
{

\renewcommand{\refname}{TÀI LIỆU THAM KHẢO}
\addcontentsline{toc}{chapter}{Tài liệu tham khảo}%section
%%%%\printbibliography
\bibliography{Tai_lieu_tham_khao/tai_lieu}
\bibliographystyle{abbrv}%{plain}%%{siam}


 \newpage
 \addcontentsline{toc}{chapter}{NHẬN XÉT CỦA GIÁO VIÊN HƯỚNG DẪN}
\begin{center}
\large {\textbf{NHẬN XÉT CỦA GIÁO VIÊN HƯỚNG DẪN}}
\end{center}

\noindent\makebox[\textwidth]{\dotfill}

\vspace{0.35cm}
\noindent\makebox[\textwidth]{\dotfill}

\vspace{0.35cm}
\noindent\makebox[\textwidth]{\dotfill}

\vspace{0.35cm}
\noindent\makebox[\textwidth]{\dotfill}

\vspace{0.35cm}
\noindent\makebox[\textwidth]{\dotfill}

\vspace{0.35cm}
\noindent\makebox[\textwidth]{\dotfill}

\vspace{0.35cm}
\noindent\makebox[\textwidth]{\dotfill}

\vspace{0.35cm}
\noindent\makebox[\textwidth]{\dotfill}

\vspace{0.35cm}
\noindent\makebox[\textwidth]{\dotfill}

\vspace{0.35cm}
\noindent\makebox[\textwidth]{\dotfill}

\vspace{0.35cm}
\noindent\makebox[\textwidth]{\dotfill}

\vspace{0.35cm}
\noindent\makebox[\textwidth]{\dotfill}

\vspace{1.2cm}

\begin{flushright}
Thái Nguyên, tháng \ldots{} năm \ldots

Giáo viên hướng dẫn

\vspace{3cm}
\end{flushright}



}

\end{document}